% --- Archon physics formalization source begin ---
% archon:physics
% archon:covers PhyXMiniProblems/problem_phyx_mini_0305.lean
% archon:source-report reports/phyx_mini/problem_phyx_mini_0305.source.json

\chapter{Physics problem phyx\_mini\_0305}
\label{ch:PhyXMiniProblems_problem_phyx_mini_0305}

\paragraph{Problem source.}
\#\# Physical scenario

A certain transverse sinusoidal wave of wavelength \$20\$ cm is moving in the positive direction of an \$x\$ axis. The transverse velocity of the particle at \$x = 0\$ as a function of time is shown in figure, where the scale of the vertical axis is set by \$u\_s = 5.0\$ cm/s.

\#\# Question

What is the amplitude?

\#\# Answer choices

- A: A: 2.0 cm

- B: B: 2.4 cm

- C: C: 2.8 cm

- D: D: 3.2 cm

\#\# Figure caption (auxiliary; use the image as primary evidence)

The image is a graph depicting a sinusoidal wave. The horizontal axis represents time (t) in seconds, labeled from 0 to 4. The vertical axis represents velocity \textbackslash{}( u \textbackslash{}) in centimeters per second (cm/s). The axis is labeled from \textbackslash{}(-u\_s\textbackslash{}) to \textbackslash{}(u\_s\textbackslash{}).

- The wave starts at the origin (0,0), rises to a peak at \textbackslash{}(u\_s\textbackslash{}) around 1 second, crosses zero at 2 seconds, reaches a trough at \textbackslash{}(-u\_s\textbackslash{}) around 3 seconds, and returns to zero at 4 seconds.
- Vertical grid lines are present at integer values of time (1, 2, 3, 4).
- The graph has a bold sine wave shape, indicating a periodic oscillation.

\#\# Dataset metadata

Wave Properties; Physical Model Grounding Reasoning, Implicit Condition Reasoning

\paragraph{Recorded answer/context.}
D: D: 3.2 cm

\paragraph{Figure/image path.}
/root/proposal\_for\_physic/science-mango/phyx\_mini\_run/phyx\_data/test\_image/305.png

\paragraph{Formalization target.}
create a compiling Lean file with sorry bodies at `PhyXMiniProblems/problem\_phyx\_mini\_0305.lean`. The Lean declarations must preserve the physical quantities, dimensions or dimensional roles, figure labels, governing-law hypotheses, and final relation expressed by this problem.
Use Mathlib/Physlib names found through LeanExplore where available. If a physics API is missing, introduce faithful local abstractions rather than scalar placeholder aliases.

\begin{theorem}[Physics formalization target]
\label{thm:physics:phyx_mini_0305:target}
\lean{PhyXMiniProblems.ProblemPhyXMini0305.problem_phyx_mini_0305}
\uses{def:physics:phyx-mini-0305:phyxminiproblems-problemphyxmini0305-lengthmagnitudereadout, def:physics:phyx-mini-0305:phyxminiproblems-problemphyxmini0305-sinusoidalstringwavesetup, def:physics:phyx-mini-0305:phyxminiproblems-problemphyxmini0305-matcheswaveproblemdata, def:physics:phyx-mini-0305:phyxminiproblems-problemphyxmini0305-matchesprimaryvelocitygraph, def:physics:phyx-mini-0305:phyxminiproblems-problemphyxmini0305-hasphysicalwaveparameters, def:physics:phyx-mini-0305:phyxminiproblems-problemphyxmini0305-satisfiessinusoidaltravelingwavelaws, def:physics:phyx-mini-0305:phyxminiproblems-problemphyxmini0305-answerchoice, def:physics:phyx-mini-0305:phyxminiproblems-problemphyxmini0305-displayedamplitudeincentimeters, def:physics:phyx-mini-0305:phyxminiproblems-problemphyxmini0305-recordeddatasetanswer, def:physics:phyx-mini-0305:phyxminiproblems-problemphyxmini0305-matchesdisplayedamplitude}
The assigned autoformalize agent should translate this physics problem into Lean declarations in the covered file, with theorem and lemma proof bodies written as `by sorry`.
\end{theorem}
\begin{proof}
This is an autoformalization task, not a proof task. Produce statements that can later be proved by the physics prover without weakening the physical model.
\end{proof}

% --- Archon declaration topology begin ---
\section{Declaration topology}
\begin{definition}[Length Magnitude]
  \label{def:physics:phyx-mini-0305:phyxminiproblems-problemphyxmini0305-lengthmagnitude}
  \lean{PhyXMiniProblems.ProblemPhyXMini0305.LengthMagnitude}
  A nonnegative physical length, used for wavelength and amplitude
\end{definition}
\begin{proof}
  This is the declared model definition of Length Magnitude.
\end{proof}
\begin{definition}[Signed Length Quantity]
  \label{def:physics:phyx-mini-0305:phyxminiproblems-problemphyxmini0305-signedlengthquantity}
  \lean{PhyXMiniProblems.ProblemPhyXMini0305.SignedLengthQuantity}
  A signed axial coordinate or transverse displacement
\end{definition}
\begin{proof}
  This is the declared model definition of Signed Length Quantity.
\end{proof}
\begin{definition}[Duration Magnitude]
  \label{def:physics:phyx-mini-0305:phyxminiproblems-problemphyxmini0305-durationmagnitude}
  \lean{PhyXMiniProblems.ProblemPhyXMini0305.DurationMagnitude}
  A nonnegative physical duration
\end{definition}
\begin{proof}
  This is the declared model definition of Duration Magnitude.
\end{proof}
\begin{definition}[Wave Number Magnitude]
  \label{def:physics:phyx-mini-0305:phyxminiproblems-problemphyxmini0305-wavenumbermagnitude}
  \lean{PhyXMiniProblems.ProblemPhyXMini0305.WaveNumberMagnitude}
  A nonnegative wave number, with radians treated as dimensionless
\end{definition}
\begin{proof}
  This is the declared model definition of Wave Number Magnitude.
\end{proof}
\begin{definition}[Angular Frequency Magnitude]
  \label{def:physics:phyx-mini-0305:phyxminiproblems-problemphyxmini0305-angularfrequencymagnitude}
  \lean{PhyXMiniProblems.ProblemPhyXMini0305.AngularFrequencyMagnitude}
  A nonnegative angular frequency, with radians treated as dimensionless
\end{definition}
\begin{proof}
  This is the declared model definition of Angular Frequency Magnitude.
\end{proof}
\begin{definition}[Speed Magnitude]
  \label{def:physics:phyx-mini-0305:phyxminiproblems-problemphyxmini0305-speedmagnitude}
  \lean{PhyXMiniProblems.ProblemPhyXMini0305.SpeedMagnitude}
  A nonnegative transverse-speed magnitude
\end{definition}
\begin{proof}
  This is the declared model definition of Speed Magnitude.
\end{proof}
\begin{definition}[Signed Velocity Quantity]
  \label{def:physics:phyx-mini-0305:phyxminiproblems-problemphyxmini0305-signedvelocityquantity}
  \lean{PhyXMiniProblems.ProblemPhyXMini0305.SignedVelocityQuantity}
  A signed transverse-velocity component
\end{definition}
\begin{proof}
  This is the declared model definition of Signed Velocity Quantity.
\end{proof}
\begin{definition}[length Magnitude Readout]
  \label{def:physics:phyx-mini-0305:phyxminiproblems-problemphyxmini0305-lengthmagnitudereadout}
  \lean{PhyXMiniProblems.ProblemPhyXMini0305.lengthMagnitudeReadout}
  \uses{def:physics:phyx-mini-0305:phyxminiproblems-problemphyxmini0305-lengthmagnitude}
  Read a nonnegative length in a selected length unit
\end{definition}
\begin{proof}
  This is the declared model definition of length Magnitude Readout.
\end{proof}
\begin{definition}[signed Length Readout]
  \label{def:physics:phyx-mini-0305:phyxminiproblems-problemphyxmini0305-signedlengthreadout}
  \lean{PhyXMiniProblems.ProblemPhyXMini0305.signedLengthReadout}
  \uses{def:physics:phyx-mini-0305:phyxminiproblems-problemphyxmini0305-signedlengthquantity}
  Read a signed coordinate or displacement in a selected length unit
\end{definition}
\begin{proof}
  This is the declared model definition of signed Length Readout.
\end{proof}
\begin{definition}[duration Readout]
  \label{def:physics:phyx-mini-0305:phyxminiproblems-problemphyxmini0305-durationreadout}
  \lean{PhyXMiniProblems.ProblemPhyXMini0305.durationReadout}
  \uses{def:physics:phyx-mini-0305:phyxminiproblems-problemphyxmini0305-durationmagnitude}
  Read a duration in a selected time unit
\end{definition}
\begin{proof}
  This is the declared model definition of duration Readout.
\end{proof}
\begin{definition}[wave Number Readout]
  \label{def:physics:phyx-mini-0305:phyxminiproblems-problemphyxmini0305-wavenumberreadout}
  \lean{PhyXMiniProblems.ProblemPhyXMini0305.waveNumberReadout}
  \uses{def:physics:phyx-mini-0305:phyxminiproblems-problemphyxmini0305-wavenumbermagnitude}
  Read wave number in radians per selected length unit
\end{definition}
\begin{proof}
  This is the declared model definition of wave Number Readout.
\end{proof}
\begin{definition}[angular Frequency Readout]
  \label{def:physics:phyx-mini-0305:phyxminiproblems-problemphyxmini0305-angularfrequencyreadout}
  \lean{PhyXMiniProblems.ProblemPhyXMini0305.angularFrequencyReadout}
  \uses{def:physics:phyx-mini-0305:phyxminiproblems-problemphyxmini0305-angularfrequencymagnitude}
  Read angular frequency in radians per selected time unit
\end{definition}
\begin{proof}
  This is the declared model definition of angular Frequency Readout.
\end{proof}
\begin{definition}[speed Readout]
  \label{def:physics:phyx-mini-0305:phyxminiproblems-problemphyxmini0305-speedreadout}
  \lean{PhyXMiniProblems.ProblemPhyXMini0305.speedReadout}
  \uses{def:physics:phyx-mini-0305:phyxminiproblems-problemphyxmini0305-speedmagnitude}
  Read a speed magnitude in coherent selected length and time units
\end{definition}
\begin{proof}
  This is the declared model definition of speed Readout.
\end{proof}
\begin{definition}[signed Velocity Readout]
  \label{def:physics:phyx-mini-0305:phyxminiproblems-problemphyxmini0305-signedvelocityreadout}
  \lean{PhyXMiniProblems.ProblemPhyXMini0305.signedVelocityReadout}
  \uses{def:physics:phyx-mini-0305:phyxminiproblems-problemphyxmini0305-signedvelocityquantity}
  Read signed transverse velocity in coherent selected units
\end{definition}
\begin{proof}
  This is the declared model definition of signed Velocity Readout.
\end{proof}
\begin{definition}[Wave Kind]
  \label{def:physics:phyx-mini-0305:phyxminiproblems-problemphyxmini0305-wavekind}
  \lean{PhyXMiniProblems.ProblemPhyXMini0305.WaveKind}
  The disturbance type stated in the problem
\end{definition}
\begin{proof}
  This is the declared model definition of Wave Kind.
\end{proof}
\begin{definition}[Propagation Direction]
  \label{def:physics:phyx-mini-0305:phyxminiproblems-problemphyxmini0305-propagationdirection}
  \lean{PhyXMiniProblems.ProblemPhyXMini0305.PropagationDirection}
  Propagation direction along the horizontal string axis
\end{definition}
\begin{proof}
  This is the declared model definition of Propagation Direction.
\end{proof}
\begin{definition}[temporal Phase Sign]
  \label{def:physics:phyx-mini-0305:phyxminiproblems-problemphyxmini0305-temporalphasesign}
  \lean{PhyXMiniProblems.ProblemPhyXMini0305.temporalPhaseSign}
  \uses{def:physics:phyx-mini-0305:phyxminiproblems-problemphyxmini0305-propagationdirection}
  Sign multiplying omega t in a traveling-wave phase
\end{definition}
\begin{proof}
  This is the declared model definition of temporal Phase Sign.
\end{proof}
\begin{definition}[Graph Axis]
  \label{def:physics:phyx-mini-0305:phyxminiproblems-problemphyxmini0305-graphaxis}
  \lean{PhyXMiniProblems.ProblemPhyXMini0305.GraphAxis}
  The two coordinate axes drawn in the supplied graph
\end{definition}
\begin{proof}
  This is the declared model definition of Graph Axis.
\end{proof}
\begin{definition}[Axis Quantity]
  \label{def:physics:phyx-mini-0305:phyxminiproblems-problemphyxmini0305-axisquantity}
  \lean{PhyXMiniProblems.ProblemPhyXMini0305.AxisQuantity}
  Physical quantity printed beside each graph axis
\end{definition}
\begin{proof}
  This is the declared model definition of Axis Quantity.
\end{proof}
\begin{definition}[Velocity Graph Landmark]
  \label{def:physics:phyx-mini-0305:phyxminiproblems-problemphyxmini0305-velocitygraphlandmark}
  \lean{PhyXMiniProblems.ProblemPhyXMini0305.VelocityGraphLandmark}
  Distinguished velocity-curve landmarks, ordered from left to right
\end{definition}
\begin{proof}
  This is the declared model definition of Velocity Graph Landmark.
\end{proof}
\begin{definition}[Transverse Velocity Time Graph]
  \label{def:physics:phyx-mini-0305:phyxminiproblems-problemphyxmini0305-transversevelocitytimegraph}
  \lean{PhyXMiniProblems.ProblemPhyXMini0305.TransverseVelocityTimeGraph}
  \uses{def:physics:phyx-mini-0305:phyxminiproblems-problemphyxmini0305-speedmagnitude, def:physics:phyx-mini-0305:phyxminiproblems-problemphyxmini0305-signedvelocityquantity, def:physics:phyx-mini-0305:phyxminiproblems-problemphyxmini0305-graphaxis, def:physics:phyx-mini-0305:phyxminiproblems-problemphyxmini0305-axisquantity, def:physics:phyx-mini-0305:phyxminiproblems-problemphyxmini0305-velocitygraphlandmark}
  The labels, units, calibration, and landmark observables carried by the primary velocity-versus-time graph. The landmark times are real scalar readouts explicitly measured in the printed unit seconds; marked velocities remain dimensionful physical quantities
\end{definition}
\begin{proof}
  This is the declared model definition of Transverse Velocity Time Graph.
\end{proof}
\begin{definition}[Sinusoidal String Wave Setup]
  \label{def:physics:phyx-mini-0305:phyxminiproblems-problemphyxmini0305-sinusoidalstringwavesetup}
  \lean{PhyXMiniProblems.ProblemPhyXMini0305.SinusoidalStringWaveSetup}
  \uses{def:physics:phyx-mini-0305:phyxminiproblems-problemphyxmini0305-lengthmagnitude, def:physics:phyx-mini-0305:phyxminiproblems-problemphyxmini0305-signedlengthquantity, def:physics:phyx-mini-0305:phyxminiproblems-problemphyxmini0305-durationmagnitude, def:physics:phyx-mini-0305:phyxminiproblems-problemphyxmini0305-wavenumbermagnitude, def:physics:phyx-mini-0305:phyxminiproblems-problemphyxmini0305-angularfrequencymagnitude, def:physics:phyx-mini-0305:phyxminiproblems-problemphyxmini0305-speedmagnitude, def:physics:phyx-mini-0305:phyxminiproblems-problemphyxmini0305-signedvelocityquantity, def:physics:phyx-mini-0305:phyxminiproblems-problemphyxmini0305-wavekind, def:physics:phyx-mini-0305:phyxminiproblems-problemphyxmini0305-propagationdirection, def:physics:phyx-mini-0305:phyxminiproblems-problemphyxmini0305-transversevelocitytimegraph}
  Independent parameters and observables of the traveling string wave. Neither the displacement amplitude nor its numerical readout is defined from the requested answer. Governing-law premises below relate these fields
\end{definition}
\begin{proof}
  This is the declared model definition of Sinusoidal String Wave Setup.
\end{proof}
\begin{definition}[Matches Wave Problem Data]
  \label{def:physics:phyx-mini-0305:phyxminiproblems-problemphyxmini0305-matcheswaveproblemdata}
  \lean{PhyXMiniProblems.ProblemPhyXMini0305.MatchesWaveProblemData}
  \uses{def:physics:phyx-mini-0305:phyxminiproblems-problemphyxmini0305-lengthmagnitudereadout, def:physics:phyx-mini-0305:phyxminiproblems-problemphyxmini0305-signedlengthreadout, def:physics:phyx-mini-0305:phyxminiproblems-problemphyxmini0305-sinusoidalstringwavesetup}
  Problem-statement data independent of the requested amplitude: the wave is a right-moving transverse sinusoid of wavelength 20 cm, observed at x = 0
\end{definition}
\begin{proof}
  This is the declared model definition of Matches Wave Problem Data.
\end{proof}
\begin{definition}[Matches Primary Velocity Graph]
  \label{def:physics:phyx-mini-0305:phyxminiproblems-problemphyxmini0305-matchesprimaryvelocitygraph}
  \lean{PhyXMiniProblems.ProblemPhyXMini0305.MatchesPrimaryVelocityGraph}
  \uses{def:physics:phyx-mini-0305:phyxminiproblems-problemphyxmini0305-durationreadout, def:physics:phyx-mini-0305:phyxminiproblems-problemphyxmini0305-speedreadout, def:physics:phyx-mini-0305:phyxminiproblems-problemphyxmini0305-signedvelocityreadout, def:physics:phyx-mini-0305:phyxminiproblems-problemphyxmini0305-velocitygraphlandmark, def:physics:phyx-mini-0305:phyxminiproblems-problemphyxmini0305-sinusoidalstringwavesetup}
  Calibrated evidence from the supplied bitmap. The four equal horizontal intervals cover one cycle from 0 s through 4 s; the curve touches +u\_s at 1 s and -u\_s at 3 s, where u\_s = 5 cm/s. These are figure/data readouts. In particular, this structure does not give the displacement amplitude or select an answer choice
\end{definition}
\begin{proof}
  This is the declared model definition of Matches Primary Velocity Graph.
\end{proof}
\begin{definition}[Has Physical Wave Parameters]
  \label{def:physics:phyx-mini-0305:phyxminiproblems-problemphyxmini0305-hasphysicalwaveparameters}
  \lean{PhyXMiniProblems.ProblemPhyXMini0305.HasPhysicalWaveParameters}
  \uses{def:physics:phyx-mini-0305:phyxminiproblems-problemphyxmini0305-lengthmagnitudereadout, def:physics:phyx-mini-0305:phyxminiproblems-problemphyxmini0305-durationreadout, def:physics:phyx-mini-0305:phyxminiproblems-problemphyxmini0305-wavenumberreadout, def:physics:phyx-mini-0305:phyxminiproblems-problemphyxmini0305-angularfrequencyreadout, def:physics:phyx-mini-0305:phyxminiproblems-problemphyxmini0305-speedreadout, def:physics:phyx-mini-0305:phyxminiproblems-problemphyxmini0305-sinusoidalstringwavesetup}
  Positivity conditions for the nondegenerate wave shown in the graph
\end{definition}
\begin{proof}
  This is the declared model definition of Has Physical Wave Parameters.
\end{proof}
\begin{definition}[Satisfies Sinusoidal Traveling Wave Laws]
  \label{def:physics:phyx-mini-0305:phyxminiproblems-problemphyxmini0305-satisfiessinusoidaltravelingwavelaws}
  \lean{PhyXMiniProblems.ProblemPhyXMini0305.SatisfiesSinusoidalTravelingWaveLaws}
  \uses{def:physics:phyx-mini-0305:phyxminiproblems-problemphyxmini0305-signedlengthquantity, def:physics:phyx-mini-0305:phyxminiproblems-problemphyxmini0305-lengthmagnitudereadout, def:physics:phyx-mini-0305:phyxminiproblems-problemphyxmini0305-signedlengthreadout, def:physics:phyx-mini-0305:phyxminiproblems-problemphyxmini0305-durationreadout, def:physics:phyx-mini-0305:phyxminiproblems-problemphyxmini0305-wavenumberreadout, def:physics:phyx-mini-0305:phyxminiproblems-problemphyxmini0305-angularfrequencyreadout, def:physics:phyx-mini-0305:phyxminiproblems-problemphyxmini0305-speedreadout, def:physics:phyx-mini-0305:phyxminiproblems-problemphyxmini0305-signedvelocityreadout, def:physics:phyx-mini-0305:phyxminiproblems-problemphyxmini0305-temporalphasesign, def:physics:phyx-mini-0305:phyxminiproblems-problemphyxmini0305-sinusoidalstringwavesetup}
  Generic governing laws for a sinusoidal traveling string wave: k lambda = 2 pi; omega T = 2 pi; y(x,t) = A sin(kx ± omega t + phi), with the minus sign for positive-x propagation; transverse particle velocity is the time derivative of displacement; the separately represented maximum transverse speed bounds the velocity magnitude and is attained. The derivative interface is grounded by Mathlib's HasDerivAt API. No field states a numerical displacement amplitude or one of the displayed answers
\end{definition}
\begin{proof}
  This is the declared model definition of Satisfies Sinusoidal Traveling Wave Laws.
\end{proof}
\begin{lemma}[angular Frequency In Radians Per Second eq pi over two]
  \label{lem:physics:phyx-mini-0305:phyxminiproblems-problemphyxmini0305-angularfrequencyinradianspersecond-eq-pi-over-two}
  \lean{PhyXMiniProblems.ProblemPhyXMini0305.angularFrequencyInRadiansPerSecond_eq_pi_over_two}
  \uses{def:physics:phyx-mini-0305:phyxminiproblems-problemphyxmini0305-angularfrequencyreadout, def:physics:phyx-mini-0305:phyxminiproblems-problemphyxmini0305-sinusoidalstringwavesetup, def:physics:phyx-mini-0305:phyxminiproblems-problemphyxmini0305-matchesprimaryvelocitygraph, def:physics:phyx-mini-0305:phyxminiproblems-problemphyxmini0305-satisfiessinusoidaltravelingwavelaws}
  The 4 s cycle gives omega = pi/2 rad/s
\end{lemma}
\begin{proof}
  The conclusion follows from the governing relations and figure readouts named in the statement of angular Frequency In Radians Per Second eq pi over two.
\end{proof}
\begin{lemma}[wave Number In Radians Per Centimeter eq pi over ten]
  \label{lem:physics:phyx-mini-0305:phyxminiproblems-problemphyxmini0305-wavenumberinradianspercentimeter-eq-pi-over-ten}
  \lean{PhyXMiniProblems.ProblemPhyXMini0305.waveNumberInRadiansPerCentimeter_eq_pi_over_ten}
  \uses{def:physics:phyx-mini-0305:phyxminiproblems-problemphyxmini0305-wavenumberreadout, def:physics:phyx-mini-0305:phyxminiproblems-problemphyxmini0305-sinusoidalstringwavesetup, def:physics:phyx-mini-0305:phyxminiproblems-problemphyxmini0305-matcheswaveproblemdata, def:physics:phyx-mini-0305:phyxminiproblems-problemphyxmini0305-satisfiessinusoidaltravelingwavelaws}
  The stated 20 cm wavelength gives k = pi/10 rad/cm. Wavelength does not enter the final amplitude calculation at a fixed string point, but this lemma records its physical role in the traveling-wave model
\end{lemma}
\begin{proof}
  The conclusion follows from the governing relations and figure readouts named in the statement of wave Number In Radians Per Centimeter eq pi over ten.
\end{proof}
\begin{lemma}[maximum Transverse Speed eq angular Frequency mul amplitude]
  \label{lem:physics:phyx-mini-0305:phyxminiproblems-problemphyxmini0305-maximumtransversespeed-eq-angularfrequency-mul-amplitude}
  \lean{PhyXMiniProblems.ProblemPhyXMini0305.maximumTransverseSpeed_eq_angularFrequency_mul_amplitude}
  \uses{def:physics:phyx-mini-0305:phyxminiproblems-problemphyxmini0305-lengthmagnitudereadout, def:physics:phyx-mini-0305:phyxminiproblems-problemphyxmini0305-angularfrequencyreadout, def:physics:phyx-mini-0305:phyxminiproblems-problemphyxmini0305-speedreadout, def:physics:phyx-mini-0305:phyxminiproblems-problemphyxmini0305-sinusoidalstringwavesetup, def:physics:phyx-mini-0305:phyxminiproblems-problemphyxmini0305-satisfiessinusoidaltravelingwavelaws}
  Differentiating the sinusoidal displacement and using that |cos| attains one shows that the transverse-speed amplitude is omega A. This is a derived lemma rather than a premise field
\end{lemma}
\begin{proof}
  The conclusion follows from the governing relations and figure readouts named in the statement of maximum Transverse Speed eq angular Frequency mul amplitude.
\end{proof}
\begin{lemma}[displacement Amplitude In Centimeters eq ten div pi]
  \label{lem:physics:phyx-mini-0305:phyxminiproblems-problemphyxmini0305-displacementamplitudeincentimeters-eq-ten-div-pi}
  \lean{PhyXMiniProblems.ProblemPhyXMini0305.displacementAmplitudeInCentimeters_eq_ten_div_pi}
  \uses{def:physics:phyx-mini-0305:phyxminiproblems-problemphyxmini0305-lengthmagnitudereadout, def:physics:phyx-mini-0305:phyxminiproblems-problemphyxmini0305-sinusoidalstringwavesetup, def:physics:phyx-mini-0305:phyxminiproblems-problemphyxmini0305-matchesprimaryvelocitygraph, def:physics:phyx-mini-0305:phyxminiproblems-problemphyxmini0305-hasphysicalwaveparameters, def:physics:phyx-mini-0305:phyxminiproblems-problemphyxmini0305-satisfiessinusoidaltravelingwavelaws}
  The graph gives maximum transverse speed 5 cm/s and period 4 s, hence omega = pi/2 rad/s. Solving u\_max = omega A gives the exact displacement amplitude A = 10/pi cm
\end{lemma}
\begin{proof}
  The conclusion follows from the governing relations and figure readouts named in the statement of displacement Amplitude In Centimeters eq ten div pi.
\end{proof}
\begin{definition}[Answer Choice]
  \label{def:physics:phyx-mini-0305:phyxminiproblems-problemphyxmini0305-answerchoice}
  \lean{PhyXMiniProblems.ProblemPhyXMini0305.AnswerChoice}
  Labels of the four amplitude choices printed with the problem
\end{definition}
\begin{proof}
  This is the declared model definition of Answer Choice.
\end{proof}
\begin{definition}[displayed Amplitude In Centimeters]
  \label{def:physics:phyx-mini-0305:phyxminiproblems-problemphyxmini0305-displayedamplitudeincentimeters}
  \lean{PhyXMiniProblems.ProblemPhyXMini0305.displayedAmplitudeInCentimeters}
  \uses{def:physics:phyx-mini-0305:phyxminiproblems-problemphyxmini0305-answerchoice}
  Displayed displacement amplitude in centimeters
\end{definition}
\begin{proof}
  This is the declared model definition of displayed Amplitude In Centimeters.
\end{proof}
\begin{definition}[recorded Dataset Answer]
  \label{def:physics:phyx-mini-0305:phyxminiproblems-problemphyxmini0305-recordeddatasetanswer}
  \lean{PhyXMiniProblems.ProblemPhyXMini0305.recordedDatasetAnswer}
  \uses{def:physics:phyx-mini-0305:phyxminiproblems-problemphyxmini0305-answerchoice}
  The answer label recorded in the supplied dataset metadata
\end{definition}
\begin{proof}
  This is the declared model definition of recorded Dataset Answer.
\end{proof}
\begin{definition}[Matches Displayed Amplitude]
  \label{def:physics:phyx-mini-0305:phyxminiproblems-problemphyxmini0305-matchesdisplayedamplitude}
  \lean{PhyXMiniProblems.ProblemPhyXMini0305.MatchesDisplayedAmplitude}
  \uses{def:physics:phyx-mini-0305:phyxminiproblems-problemphyxmini0305-answerchoice, def:physics:phyx-mini-0305:phyxminiproblems-problemphyxmini0305-displayedamplitudeincentimeters}
  Agreement with an amplitude displayed to the nearest tenth of a centimeter
\end{definition}
\begin{proof}
  This is the declared model definition of Matches Displayed Amplitude.
\end{proof}
% --- Archon declaration topology end ---
% --- Archon physics formalization source end ---
